% Text converted from the supplied Report 90072 .pdf.
\chapter{Introduction and data}\label{ch:1}

\section{Project question}

An ECG records the electrical activity of the heart. Its R peaks provide a clear marker for the timing of each beat. Breathing can slightly change the time between beats, an effect called respiratory sinus arrhythmia \cite{projectnote}. Breathing movements can also change the size of the ECG waveform \cite{varon2015}. These changes may contain information about breathing frequency.

The project asks whether the strongest frequency in either of these ECG features agrees with an independently recorded respiratory signal. The main tasks are to read the data, identify the beat features, calculate their frequency spectra and compare the results. A spectrum shows how strongly different frequencies contribute to a signal.

\section{Dataset}

Record \texttt{a01} comes from the Apnea-ECG database \cite{penzel2000,physionet2000}. The ECG values are stored in \texttt{a01.dat}, with recording information in a01.hea. The file \texttt{a01.qrs} contains supplied heartbeat marks, while \texttt{a01.apn} contains one apnea label for each annotated minute.

\begin{table}[H]
\centering
\caption{Main properties of the recording.}
\small
\begin{tabularx}{\linewidth}{@{}p{0.39\linewidth}X@{}}
\toprule
\textbf{Property} & \textbf{Value} \\
\midrule
ECG channels and sampling rate & One channel at 100 Hz \\
Full ECG length & 2,957,000 samples; 8 h 12 min 50 s \\
Segment analysed for respiration & First 60 minutes; 30 two-minute windows \\
Supplied heartbeat marks & 29,938 across the full ECG \\
Minute labels across the full record & 470 A labels and 19 N labels \\
Primary respiratory reference & Oronasal airflow from the paired \texttt{a01r} record \\
\bottomrule
\end{tabularx}
\end{table}

The paired \texttt{a01r} record also includes chest effort, abdominal effort and oxygen saturation. Airflow was used as the main respiratory reference; oxygen saturation was not used for frequency estimation. Both ECG and respiratory records cover the first hour.

\section{Scope}

The first-hour limit was a practical choice because downloading the full respiratory file was slow. The first two-minute window was used while developing the method, and the same settings were then used for the remaining windows. This is an exploratory study of one recording, rather than a test of performance across different people.

